% Does GARCH(1,1) Actually Forecast Volatility? A Controlled Study.
% Controlled, seeded, synthetic-DGP study of GARCH volatility forecasting.
% Every number in this manuscript is verified against results/results.json by
% scripts/check_paper_numbers.py; do not edit numbers without re-running it.
\documentclass[11pt]{article}
\usepackage[margin=1in]{geometry}
\usepackage{amsmath,amssymb}
\usepackage{booktabs}
\usepackage[round]{natbib}
\usepackage{microtype}
\usepackage[hidelinks]{hyperref}
\setlength{\emergencystretch}{2em}

\newcommand{\E}{\mathbb{E}}
\newcommand{\Var}{\mathrm{Var}}
\newcommand{\lrv}{\bar\sigma^{2}}
\newcommand{\pers}{(\alpha+\beta)}

\title{Does GARCH(1,1) Actually Forecast Volatility?\\
\large A Controlled Study with Known Ground Truth}
\author{Eugen Soloviov}
\date{July 2026}

\begin{document}
\maketitle

\begin{abstract}
GARCH(1,1) is the workhorse of applied volatility modeling, but ``it forecasts
well'' is usually asserted on real data where the true conditional variance is
never observed. We instead run four seeded, fully reproducible experiments on a
synthetic GARCH(1,1) data-generating process (DGP) whose conditional-variance
path is known by construction, and calibrate exactly when the model earns its
keep and whether the evaluation can be trusted without observing volatility.
First, maximum-likelihood estimation is consistent: on a crypto-like DGP
($\alpha=0.09$, $\beta=0.90$, persistence $0.99$) the persistence estimate
converges to the truth ($0.977 \to 0.989$ as the sample grows from $500$ to
$5000$) and every parameter's root-mean-square error falls roughly like
$1/\sqrt{T}$ (persistence RMSE $0.028 \to 0.004$). Second, in a one-step-ahead
forecast contest scored against the \emph{true} conditional variance by the
proxy-robust QLIKE loss, the correctly specified GARCH wins at every DGP,
beating the best naive competitor (RiskMetrics EWMA or a rolling window) by
$83$--$93\%$; the advantage is largest at moderate persistence and shrinks
toward the near-random regime (best-naive QLIKE $0.0222 \to 0.0079$ as
persistence falls from $0.90$ to $0.05$), and near-IGARCH the RiskMetrics EWMA
is the closest competitor because it \emph{is} a persistence-one GARCH. Third,
across a persistence grid the persistence itself is pinned down tightly even
near the unit root (RMSE $0.032 \to 0.005$ as persistence rises to $0.995$),
but everything derived from it explodes: the implied volatility half-life runs
from $6.6$ to $138.3$ observations while its estimate becomes wildly dispersed
(interquartile range $2.6 \to 121.4$) and the long-run variance grows
increasingly downward biased --- the near-IGARCH fragility, quantified. Fourth,
re-scoring the same forecasters against the noisy squared-return proxy instead
of the true variance leaves the ranking intact for the Patton-robust losses:
the proxy inflates every QLIKE by a nearly forecaster-independent constant (at
the crypto DGP, GARCH's QLIKE moves $0.0018 \to 1.269$ and EWMA's
$0.0119 \to 1.280$, preserving the gap) and GARCH is selected under the proxy
at all six DGPs, whereas the non-robust MAE agrees on only $82\%$ of noisy
draws versus $96\%$ for QLIKE. GARCH's win here is guaranteed by construction
--- the DGP is GARCH --- so the contribution is not that GARCH beats naive
forecasters on real crypto, but the calibration of \emph{when} and \emph{by how
much} it should, and the demonstration that the evaluation is trustworthy even
when true volatility is unobservable. This study accompanies a
\href{https://marketmaker.cc}{marketmaker.cc} blog post.
\end{abstract}

\section{Introduction}
\label{sec:intro}

Financial returns are close to unpredictable in sign but strongly predictable
in magnitude: large moves cluster, small moves cluster, and the squared returns
carry persistent autocorrelation that the raw returns do not
\citep{mandelbrot1963variation}. The GARCH(1,1) model of
\citet{bollerslev1986garch}, generalizing the ARCH model of
\citet{engle1982arch}, captures this with three parameters, and has become the
default conditional-variance model in risk management. In cryptocurrency
markets the fit is if anything sharper, and typically lands in the near-IGARCH
regime with estimated persistence close to one
\citep{katsiampa2017volatility,chu2017garch}.

The standard evidence that GARCH ``forecasts volatility well''
\citep{andersenbollerslev1998answering,hansenlunde2005garch} is collected on
real data, where the object being forecast --- the conditional variance
$\sigma_t^2$ --- is latent and never observed. Two questions are therefore hard
to answer cleanly on real markets. When exactly does the model beat a naive
volatility estimate, and by how much? And can we even trust a forecast
comparison when the target is unobservable and must be replaced by a noisy
proxy such as the squared return? This paper answers both under controlled
conditions, by generating the data from a known GARCH(1,1) so that the true
conditional-variance path is available as ground truth.

We make no claim about real crypto. On the contrary: because our DGP \emph{is}
GARCH, the correctly specified model is guaranteed to win the forecast contest,
and reporting that it wins would be circular. The contribution is the
calibration around that fact --- the size of the advantage and how it varies
with persistence, the consistency of the estimator, the near-IGARCH fragility
of every persistence-derived quantity, and, most usefully, a controlled
confirmation of the \citet{patton2011volatility} result that the QLIKE and MSE
losses rank forecasters identically whether they are scored against the true
variance or against the noisy squared-return proxy. That last result is what
licenses trusting a volatility-forecast evaluation on real data at all.

Concretely, we run four experiments (Section~\ref{sec:design}) on the
estimators and forecasters of Section~\ref{sec:methods}:
\begin{enumerate}
\item \textbf{Parameter recovery.} Simulate a crypto-like GARCH(1,1) and refit
  by maximum likelihood over $300$ seeds at sample sizes
  $T \in \{500, 1000, 2000, 5000\}$. Bias and RMSE shrink with $T$ and the
  persistence estimate converges to the true $0.99$: MLE consistency, measured.
\item \textbf{Forecast contest.} Score one-step-ahead forecasters --- the
  correctly specified GARCH refit, RiskMetrics EWMA ($\lambda = 0.94$), and
  rolling-window sample variances --- against the true conditional variance
  with MSE and QLIKE, across DGPs from near-IGARCH down to near-random.
  Quantify the GARCH advantage, its dependence on persistence, and the analytic
  multi-step mean-reversion of the forecast toward the long-run variance.
\item \textbf{Persistence and half-life.} Across a grid of true persistence
  ($0.90$ to $0.995$), report the estimated persistence, the implied half-life
  $\ln 0.5 / \ln\pers$, and how estimation and long-horizon uncertainty blow up
  as persistence approaches one.
\item \textbf{Proxy robustness.} Re-score the contest against the squared-return
  proxy and confirm (Patton 2011) that the robust losses preserve the ranking,
  so the evaluation is trustworthy without observing true volatility.
\end{enumerate}
All results derive from one seeded run of one public harness
(\texttt{scripts/run\_all.py}); a companion script
(\texttt{scripts/check\_paper\_numbers.py}) verifies every numeric claim in this
manuscript against the run's JSON output.

\section{Model, estimator, and forecasters}
\label{sec:methods}

\paragraph{The data-generating process.} We simulate GARCH(1,1) with zero
conditional mean and Gaussian innovations,
\begin{equation}
\label{eq:garch}
r_t = \sigma_t z_t, \quad z_t \sim \text{iid } \mathcal{N}(0,1), \qquad
\sigma_t^2 = \omega + \alpha\, r_{t-1}^2 + \beta\, \sigma_{t-1}^2,
\end{equation}
with $\omega > 0$, $\alpha,\beta \ge 0$, and $\alpha + \beta < 1$. Because we
generate the path, the conditional variance $\sigma_t^2$ used to draw each
$r_t$ is retained as ground truth. Taking unconditional expectations of
Eq.~\eqref{eq:garch} gives the long-run variance, and iterating gives the
persistence and half-life,
\begin{equation}
\label{eq:props}
\lrv = \frac{\omega}{1-\alpha-\beta}, \qquad
\text{persistence} = \alpha+\beta, \qquad
h_{1/2} = \frac{\ln 0.5}{\ln\pers},
\end{equation}
the half-life being the number of observations for a variance shock to decay
halfway back to $\lrv$. Our base ``crypto-like'' DGP fixes $\omega = 0.04$,
$\alpha = 0.09$, $\beta = 0.90$, so the persistence is $0.99$ (near-IGARCH),
the long-run variance is $4.0$ (per-observation volatility $2.0$), and the true
half-life is $69.0$ observations (derived from Eq.~\eqref{eq:props}). A
burn-in is discarded so each retained path starts from the stationary
distribution.

\paragraph{Estimation.} Parameters are estimated by Gaussian maximum
likelihood. With $r_t \mid \mathcal{F}_{t-1} \sim \mathcal{N}(0,
\sigma_t^2(\theta))$ the log-likelihood is a sum of one-step densities,
$\ell(\theta) = -\tfrac12 \sum_t [\ln 2\pi + \ln\sigma_t^2(\theta) + r_t^2 /
\sigma_t^2(\theta)]$, maximized numerically by the \texttt{arch} library of
\citet{sheppard2023arch}. The DGP is scaled so returns are $O(1)$, so the
$\times 100$ rescaling that daily returns usually require is unnecessary and is
switched off.

\paragraph{Forecasters.} Each forecaster produces a causal one-step-ahead
variance $h_t = \widehat{\Var}(r_t \mid \mathcal{F}_{t-1})$:
\begin{itemize}
\item \emph{GARCH} (correctly specified): fit on a training block, then filter
  $h_t = \hat\omega + \hat\alpha\, r_{t-1}^2 + \hat\beta\, h_{t-1}$ over the
  test block. With the true parameters this reproduces $\sigma_t^2$ exactly
  after the seed transient; with estimated parameters it carries only
  estimation error.
\item \emph{EWMA / RiskMetrics} \citep{riskmetrics1996}, $\lambda = 0.94$:
  $h_t = (1-\lambda) r_{t-1}^2 + \lambda\, h_{t-1}$. This is the constrained
  GARCH special case $\omega = 0$, $\alpha = 1-\lambda$, $\beta = \lambda$, with
  persistence exactly one.
\item \emph{Rolling window}, $w \in \{22, 66, 132\}$ observations (about one,
  three, and six trading months): $h_t = \text{mean}(r_{t-w}^2, \dots,
  r_{t-1}^2)$.
\end{itemize}

\paragraph{Multi-step forecast.} Iterating Eq.~\eqref{eq:garch} in conditional
expectation, using $\E_T[r_{T+k}^2] = \E_T[\sigma_{T+k}^2]$, gives the analytic
term structure
\begin{equation}
\label{eq:multistep}
\E_T[\sigma_{T+h}^2] = \lrv + \pers^{\,h-1}\bigl(\sigma_{T+1}^2 - \lrv\bigr),
\end{equation}
the forecast decaying geometrically from today's one-step variance toward the
long-run level at rate $\pers^{\,h-1}$.

\paragraph{Losses.} A forecast $h$ is scored against a target $X$ by
\citep{patton2011volatility}
\begin{equation}
\label{eq:losses}
\mathrm{MSE}(X,h) = \E[(X-h)^2], \quad
\mathrm{QLIKE}(X,h) = \E\!\left[\frac{X}{h} - \ln\frac{X}{h} - 1\right], \quad
\mathrm{MAE}(X,h) = \E[|X-h|].
\end{equation}
QLIKE is the Bregman divergence generated by $-\ln$: non-negative, uniquely
minimized at $h = X$, and dependent only on the ratio $X/h$ (scale free). The
target $X$ is the true conditional variance $\sigma_t^2$ in Experiments~2--3
and the squared-return proxy $r_t^2$ in Experiment~4. MSE and QLIKE are the
Patton-robust losses (ranking invariant to a conditionally unbiased proxy);
MAE, which targets the conditional median, is not and is carried as a contrast.

\section{Experimental design}
\label{sec:design}

All experiments are generated and analyzed by a single seeded Python harness
(\texttt{scripts/run\_all.py}, methods in \texttt{scripts/garch.py}) under
Python 3.14.6 with NumPy 2.5.1 and \texttt{arch} 8.0.0, so every number is
bit-reproducible by one command. \emph{Parameter recovery} refits the base DGP
over $M = 300$ seeds at each $T \in \{500, 1000, 2000, 5000\}$, reporting the
bias, RMSE, and standard deviation of each estimate. \emph{Forecast contest}
simulates $M = 200$ paths of length $3000$ per DGP, fits GARCH on the first
$2000$ observations, and scores all forecasters on the final $1000$ (so
GARCH's parameters never see the test block); it sweeps DGPs of persistence
$\{0.99, 0.90, 0.70, 0.40, 0.20, 0.05\}$, holding the long-run variance fixed
at $4.0$ and $\alpha$'s share of persistence fixed at the base ratio. The
multi-step sub-experiment forecasts $h \in \{1, 5, 10, 22\}$ steps ahead from
every test origin on the crypto-like DGP over $M = 120$ paths.
\emph{Persistence and half-life} refits over $M = 300$ seeds of length $3000$
across true persistence $\{0.90, 0.95, 0.97, 0.98, 0.99, 0.995\}$.
\emph{Proxy robustness} re-scores the contest against $r_t^2$ and compares
forecaster orderings under the two targets.

\section{Results}
\label{sec:results}

\subsection{Parameter recovery: MLE is consistent}
\label{sec:res-recovery}

Table~\ref{tab:recovery} reports the estimation error on the crypto-like DGP.
Every parameter is recovered with a bias and RMSE that shrink steadily with the
sample size, at roughly the $1/\sqrt{T}$ rate a consistent estimator should
show: the persistence RMSE falls from $0.028$ at $T = 500$ to $0.004$ at
$T = 5000$, and the persistence estimate itself climbs from $0.977$ toward the
true $0.99$ (reaching $0.989$), with a small negative small-sample bias
($-0.013$ at $T = 500$) that is the well-known downward bias of the persistence
estimate near the unit root. The reaction and memory parameters land on their
targets ($\hat\alpha \to 0.089$, $\hat\beta \to 0.900$ at $T = 5000$), and the
intercept $\omega$, which is the hardest to pin down in the near-IGARCH regime,
carries the largest relative error (RMSE $0.078 \to 0.011$, a persistent small
upward bias $0.036 \to 0.003$). The implied half-life, a strongly nonlinear
function of the estimates, is correspondingly slow to converge: its median
climbs from $41.3$ toward the true $69.0$, reaching only $63.0$ at $T = 5000$
--- a first sign of the near-IGARCH fragility that Section~\ref{sec:res-persist}
quantifies.

\begin{table}[t]
\centering
\caption{Experiment 1 (parameter recovery): bias and root-mean-square error of
the Gaussian-MLE estimates on the crypto-like DGP ($\omega = 0.04$,
$\alpha = 0.09$, $\beta = 0.90$, persistence $0.99$, true half-life $69.0$
observations), over $M = 300$ seeds at each sample size. Errors fall with $T$;
the persistence estimate converges to $0.99$.}
\label{tab:recovery}
\begin{tabular}{rrrrrrr}
\toprule
$T$ & $\widehat{\text{pers.}}$ & pers.\ RMSE & $\alpha$ RMSE & $\beta$ RMSE
    & $\omega$ RMSE & median $h_{1/2}$ \\
\midrule
500  & 0.977 & 0.028  & 0.031  & 0.040  & 0.078 & 41.3 \\
1000 & 0.984 & 0.013  & 0.018  & 0.022  & 0.031 & 50.7 \\
2000 & 0.987 & 0.0072 & 0.014  & 0.014  & 0.017 & 56.8 \\
5000 & 0.989 & 0.0044 & 0.0083 & 0.0092 & 0.011 & 63.0 \\
\bottomrule
\end{tabular}
\end{table}

\subsection{Forecast contest: GARCH wins, and by how much}
\label{sec:res-contest}

Table~\ref{tab:contest} scores the one-step-ahead forecasters against the true
conditional variance. The correctly specified GARCH minimizes both QLIKE and
MSE at every DGP --- as it must, since it is the DGP --- but the interesting
content is the margin and its shape. On the crypto-like DGP (persistence
$0.99$) GARCH's QLIKE is $0.0018$ against $0.0119$ for the best naive
competitor, an $85.0\%$ reduction; on MSE the gap is far starker ($0.219$
versus $1.13$ for EWMA), because MSE on the variance level is dominated by the
rare large-variance spikes that only a conditional model tracks.

The margin is not monotone in persistence, and the reason is instructive. It
peaks at moderate persistence --- $92.8\%$ at persistence $0.90$, where the
best naive QLIKE is worst at $0.0222$ --- and falls off in both directions.
Toward the near-random end the absolute advantage collapses because there is
little volatility dynamics left to model: the best-naive QLIKE falls
monotonically from $0.0222$ at persistence $0.90$ to $0.0079$ at persistence
$0.05$, and GARCH's edge over it shrinks with it (absolute QLIKE gap $0.0206
\to 0.0065$). Toward the near-IGARCH end the gap also narrows, to $85.0\%$ at
persistence $0.99$, but for the opposite reason: there the best naive is the
RiskMetrics EWMA, and near a unit root EWMA --- a persistence-one GARCH --- is
a good approximation, so it trails GARCH closely rather than badly. Below
persistence $0.70$ the best naive is instead the longest rolling window, a
near-constant-variance estimate that suits a nearly-static DGP.

\begin{table}[t]
\centering
\caption{Experiment 2 (forecast contest): one-step-ahead forecast losses
against the TRUE conditional variance, averaged over $M = 200$ paths, by DGP
persistence (long-run variance held at $4.0$). ``Best naive'' is the
lowest-QLIKE non-GARCH forecaster and its QLIKE; the gap is the percentage
QLIKE reduction of GARCH over it. GARCH also minimizes MSE at every row.}
\label{tab:contest}
\begin{tabular}{rrllrr}
\toprule
Persistence & GARCH QLIKE & Best naive & Best-naive QLIKE & QLIKE gap
            & GARCH MSE \\
\midrule
0.99 & 0.0018 & EWMA          & 0.0119 & 85.0\% & 0.219 \\
0.90 & 0.0016 & EWMA          & 0.0222 & 92.8\% & 0.070 \\
0.70 & 0.0016 & rolling (132) & 0.0148 & 89.2\% & 0.060 \\
0.40 & 0.0012 & rolling (132) & 0.0093 & 87.1\% & 0.043 \\
0.20 & 0.0011 & rolling (132) & 0.0085 & 87.3\% & 0.037 \\
0.05 & 0.0013 & rolling (132) & 0.0079 & 82.9\% & 0.042 \\
\bottomrule
\end{tabular}
\end{table}

\paragraph{Multi-step forecasts mean-revert.} Table~\ref{tab:multistep}
reports the analytic multi-step forecast of Eq.~\eqref{eq:multistep} on the
crypto-like DGP. The forecast error against the realized true variance grows
with the horizon, from an MSE of $0.58$ at $h = 1$ to $12.65$ at $h = 22$, as
the intervening shocks accumulate. Meanwhile the forecast decays toward the
long-run variance $4.0$ at the predicted geometric rate $\pers^{\,h-1}$: with
persistence $0.99$ that factor is $0.961$ at $h = 5$, $0.914$ at $h = 10$, and
$0.810$ at $h = 22$, and the empirical mean deviation of the forecast from
$\lrv$ tracks it closely (falling to $0.767$ of its one-step value by
$h = 22$). This is the near-IGARCH signature: even three weeks out, the
forecast has reverted less than a fifth of the way to its long-run level, so a
short-horizon crypto GARCH forecast is essentially ``today's variance, very
slowly fading.''

\begin{table}[t]
\centering
\caption{Multi-step forecast on the crypto-like DGP ($M = 120$ paths). MSE is
against the realized true variance $h$ steps ahead; the analytic decay is
$\pers^{\,h-1}$ with persistence $0.99$; the deviation ratio is the mean
$|\text{forecast} - \lrv|$ relative to its one-step value, confirming the
geometric mean-reversion toward the long-run variance $4.0$.}
\label{tab:multistep}
\begin{tabular}{rrrr}
\toprule
Horizon $h$ & Forecast MSE & Analytic decay $\pers^{h-1}$
            & Deviation ratio \\
\midrule
1  & 0.58  & 1.000 & 1.000 \\
5  & 3.22  & 0.961 & 0.947 \\
10 & 6.40  & 0.914 & 0.887 \\
22 & 12.65 & 0.810 & 0.767 \\
\bottomrule
\end{tabular}
\end{table}

\subsection{Persistence and half-life: the near-IGARCH fragility}
\label{sec:res-persist}

Table~\ref{tab:persist} sweeps the true persistence and exposes the crypto
practitioner's central hazard. The persistence \emph{itself} is estimated
tightly, and in fact more tightly as it approaches one: its RMSE falls from
$0.032$ at persistence $0.90$ to $0.005$ at $0.995$, because a near-unit-root
process pins the estimate against the stationarity boundary. But every quantity
\emph{derived} from persistence through the $1/(1-\pers)$ factor becomes
treacherous. The volatility half-life $\ln 0.5 / \ln\pers$ runs from $6.6$
observations at persistence $0.90$ to $138.3$ at $0.995$, and its estimate is
both downward biased (median $105.5$ against the true $138.3$ at persistence
$0.995$) and wildly dispersed: the interquartile range of the estimated
half-life explodes from $2.6$ observations at persistence $0.90$ to $121.4$ at
$0.995$ --- nearly as large as the quantity itself. The long-run variance,
whose sensitivity to persistence is $\lrv/(1-\pers)$ and thus grows from $40$
to $800$ across the grid, drifts increasingly below its true value of $4.0$
(median $3.56$ at persistence $0.995$). The lesson matches the folk warning: in
the near-IGARCH regime the point estimate of long-run volatility and half-life
should never be trusted without an interval, even though the persistence that
generates them looks precise.

\begin{table}[t]
\centering
\caption{Experiment 3 (persistence and half-life): estimation across the true
persistence grid ($M = 300$ seeds, $T = 3000$, long-run variance $4.0$). The
persistence is estimated ever more tightly toward the unit root, but the
implied half-life $\ln 0.5/\ln\pers$ and long-run variance become severely
dispersed and biased --- the near-IGARCH fragility.}
\label{tab:persist}
\begin{tabular}{rrrrrrr}
\toprule
Persistence & $\widehat{\text{pers.}}$ & pers.\ RMSE & true $h_{1/2}$
            & median $\hat h_{1/2}$ & $\hat h_{1/2}$ IQR & median $\widehat{\lrv}$ \\
\midrule
0.900 & 0.896 & 0.032  & 6.6   & 6.7   & 2.6   & 3.97 \\
0.950 & 0.947 & 0.015  & 13.5  & 13.4  & 5.1   & 4.01 \\
0.970 & 0.968 & 0.011  & 22.8  & 21.8  & 10.1  & 3.97 \\
0.980 & 0.978 & 0.0080 & 34.3  & 32.4  & 14.1  & 3.90 \\
0.990 & 0.988 & 0.0061 & 69.0  & 62.4  & 34.9  & 3.79 \\
0.995 & 0.993 & 0.0053 & 138.3 & 105.5 & 121.4 & 3.56 \\
\bottomrule
\end{tabular}
\end{table}

\subsection{Proxy robustness: the evaluation survives an unobservable target}
\label{sec:res-proxy}

Real volatility is never observed, so any forecast comparison on live data must
score against a proxy. The squared return $r_t^2$ is a conditionally unbiased
but extremely noisy one --- $r_t^2 = \sigma_t^2 z_t^2$ with
$\E[z_t^2 \mid \mathcal{F}_{t-1}] = 1$ --- and Experiment~4 asks whether it
still ranks the forecasters correctly. Table~\ref{tab:proxy} re-scores the
crypto-like contest against $r_t^2$. The effect of the proxy is dramatic in
level but not in order: it inflates every QLIKE by a nearly
forecaster-independent constant of about $1.267$ (the mean of the proxy's own
noise entropy), moving GARCH from $0.0018$ to $1.269$ and EWMA from $0.0119$ to
$1.279$, so the gap between them barely changes ($0.0101$ against true
variance, $0.0108$ against the proxy). The entire ordering is preserved, and
GARCH remains the selected forecaster; across all six DGPs GARCH is chosen
under the proxy exactly as under the true variance ($6$ of $6$), for both robust
losses.

The robustness is a property of the loss, not luck. Repeating the comparison
seed by seed on the noisy proxy --- where finite-sample noise can flip the
ordering of close competitors --- the Patton-robust losses agree with the
true-variance ordering on $96.2\%$ (QLIKE) and $95.2\%$ (MSE) of the $12000$
forecaster-pair comparisons, while the non-robust MAE agrees on only $81.5\%$.
The correctly specified model is dominant enough here that even MAE picks it as
the overall winner, but MAE is measurably less reliable on the pairwise
orderings, exactly as its lack of proxy-robustness predicts. The practical
consequence is the one that makes volatility-forecast evaluation possible at
all: scored with QLIKE or MSE, a forecast comparison run on observable squared
returns reaches the same verdict it would reach against the latent variance
nobody can see.

\begin{table}[t]
\centering
\caption{Experiment 4 (proxy robustness), crypto-like DGP: one-step-ahead
QLIKE scored against the true conditional variance versus the squared-return
proxy. The proxy adds a nearly forecaster-independent offset ($\approx 1.267$)
but preserves the ordering; GARCH stays first. Across all six DGPs GARCH is the
QLIKE-selected forecaster under the proxy as under the truth.}
\label{tab:proxy}
\begin{tabular}{lrr}
\toprule
Forecaster & QLIKE vs.\ true variance & QLIKE vs.\ $r_t^2$ proxy \\
\midrule
GARCH          & 0.0018 & 1.269 \\
EWMA           & 0.0119 & 1.279 \\
rolling (22)   & 0.0290 & 1.297 \\
rolling (66)   & 0.0738 & 1.342 \\
rolling (132)  & 0.1365 & 1.407 \\
\bottomrule
\end{tabular}
\end{table}

\section{Discussion}
\label{sec:discussion}

\paragraph{What GARCH does and does not buy you.} Under a GARCH DGP the
correctly specified model wins the one-step contest at every persistence, but
the margin over a good naive estimator is not uniform. It is largest where the
volatility is genuinely dynamic yet the naive tools are misspecified (moderate
persistence), and it shrinks at both extremes: toward a static DGP because
there is little to forecast, and toward the near-IGARCH regime because there the
RiskMetrics EWMA is itself nearly a GARCH and closes most of the gap. This is
the honest calibration behind the workhorse status of GARCH(1,1): the model
earns the most where the data have exploitable but non-trivial dynamics, and in
the near-unit-root regime typical of crypto a well-tuned EWMA is a respectable
one-parameter stand-in for one-step forecasting --- consistent with the classic
finding that little beats GARCH(1,1) but that the simplest conditional models
already capture most of the forecastable variance
\citep{hansenlunde2005garch,andersenbollerslev1998answering}.

\paragraph{Persistence is easy; everything downstream is not.} The
near-IGARCH results sharpen a warning that is easy to state and easy to ignore.
The persistence is the best-identified feature of a high-persistence GARCH fit,
tightening as it approaches one. But the half-life and the long-run variance
are $1/(1-\pers)$-type functionals, and a persistence known to $\pm 0.005$
still leaves the half-life uncertain by tens of observations and the long-run
variance uncertain by a large factor. Multi-step forecasts and any
volatility-targeting rule that leans on the long-run level inherit that
fragility; the safe move is to treat the near-IGARCH long-run variance as an
interval, refit on rolling windows, and be alert that apparent near-integration
can be an artifact of structural breaks rather than genuine persistence.

\paragraph{Why the evaluation can be trusted.} The most portable result is the
proxy-robustness of Experiment~4. It is what allows the calibration in this
paper --- all of it obtained against a latent variance we happen to know
because we simulated it --- to transfer to real data, where the same QLIKE or
MSE comparison can be run against observable squared returns and will rank
forecasters the same way. Using a non-robust loss such as MAE forfeits this
guarantee; our seed-by-seed comparison shows it degrading measurably even in a
setting where the correct model is otherwise dominant.

\section{Limitations}
\label{sec:limitations}

\begin{itemize}
\item \textbf{Synthetic GARCH DGP, favoring GARCH by construction.} The DGP is
  a Gaussian GARCH(1,1), so the correctly specified model is guaranteed to win
  the forecast contest; we never claim GARCH beats naive forecasters on real
  crypto. The contribution is the calibration of the margin and its shape, the
  consistency and near-IGARCH fragility results, and the proxy-robustness
  demonstration --- all of which require known ground truth and none of which
  is a market claim.
\item \textbf{Gaussian innovations, symmetric variance.} Real crypto returns
  have heavy tails and an asymmetric (leverage) response that plain Gaussian
  GARCH(1,1) does not model; those are deliberately out of scope here and are
  the subject of the companion work on Student-$t$/skewed innovations and
  GJR/EGARCH \citep{nelson1991egarch}. Our experiments say nothing about model \emph{misspecification},
  only about behavior under correct specification.
\item \textbf{One base parameterization per experiment.} The forecast contest
  varies persistence but holds the long-run variance and the $\alpha$-share
  fixed; the recovery and half-life grids fix the remaining parameters. We did
  not sweep the full parameter space, the innovation distribution, or the
  train/test split, so the quantitative margins are specific to these
  configurations.
\item \textbf{One-step focus; forecasters are fixed-parameter.} The contest is
  primarily one-step-ahead, and GARCH is fit once on the training block rather
  than refit at every step; a fully rolling refit and a broader multi-horizon
  loss study would sharpen the multi-step numbers.
\item \textbf{Proxy robustness shown for one proxy.} We use the squared return,
  the canonical conditionally-unbiased daily proxy; realized-variance proxies
  from intraday data are more precise and would compress the level offset of
  Experiment~4, though Patton's ranking-robustness holds for any conditionally
  unbiased proxy.
\end{itemize}

\section{Conclusion}
\label{sec:conclusion}

We asked whether GARCH(1,1) actually forecasts volatility, and answered it
where the answer can be checked: on a synthetic GARCH DGP with a known
conditional-variance path. Maximum likelihood recovers the parameters
consistently, the persistence estimate converging to the true $0.99$ as the
RMSE falls like $1/\sqrt{T}$. Scored against the true variance by the
proxy-robust QLIKE, the correctly specified model beats the best naive
forecaster by $83$--$93\%$, with the margin largest at moderate persistence and
narrowing toward both the static and the near-IGARCH extremes --- and near the
unit root a RiskMetrics EWMA is its closest competitor because it is a
persistence-one GARCH. The persistence is pinned down tightly even near
integration, but the half-life and long-run variance derived from it become
severely dispersed and biased, quantifying the near-IGARCH fragility that
should make any crypto volatility-targeting rule cautious. Finally, re-scoring
against the noisy squared-return proxy leaves the QLIKE and MSE rankings intact
--- the proxy shifts every loss by a near-constant and preserves the order ---
so the evaluation is trustworthy even though true volatility is never observed.
GARCH's victory here is built in; what is not built in, and what transfers, is
the calibration of when it matters and the confirmation that we can measure it
without seeing the target.

\paragraph{Reproducibility.} All code, tests, and outputs accompany this paper:
\texttt{scripts/run\_all.py} regenerates \texttt{results/results.json} from
fixed seeds (Python 3.14.6, NumPy 2.5.1, \texttt{arch} 8.0.0);
\texttt{scripts/check\_paper\_numbers.py} verifies every numeric claim in this
manuscript against that file and fails on any mismatch; \texttt{tests/}
contains deterministic invariant tests for the DGP, estimator, forecasters,
and losses.

\bibliographystyle{plainnat}
\bibliography{refs}

\end{document}
